\begin{abcliste}
\abc The reversed battery counts negative: it cancels one of the others.
\begin{align}
 \Delta V &= \Delta V_0 + \Delta V_0 - \Delta V_0 = \result{\Delta V_0}
\end{align}

\abc A is connected straight across the batteries: it gets the whole $\Delta V_0$, just like the reference bulb, and so the same current and the same power $P = \Delta V\cdot I$. A is \result{\text{as bright as } L_0}.

The branch with B and C has the same voltage $\Delta V_0$. B and C are identical and in series, so they share it equally: $\Delta V_B = \Delta V_C = \frac{1}{2}\,\Delta V_0$. With less voltage they also carry less current than $L_0$, so both are \result{\text{less bright than } L_0}, and equally bright. The ranking:
\begin{align}
 \result{P_A = P_0 > P_B = P_C}
\end{align}
(The branch B--C does not change the brightness of A: an ideal battery keeps its voltage, whatever is connected in parallel.)

\abc The student counts all three batteries. But the top battery is reversed: across it the potential drops by $\Delta V_0$ where the others raise it, so the three give only $\Delta V_0$, not $3\,\Delta V_0$. A is as bright as $L_0$.

Besides, even three times the voltage would not mean three times the brightness: with more voltage the current is larger too, and the power $P = \Delta V\cdot I$ grows with both (for a constant resistance by a factor of 9).
\end{abcliste}
